% This is a variation of the NeurIPS'24 format

\documentclass{article}
\PassOptionsToPackage{numbers,square,sort&compress}{natbib}
\usepackage[final]{adrl}

\usepackage[utf8]{inputenc} % allow utf-8 input
\usepackage[T1]{fontenc}    % use 8-bit T1 fonts
\usepackage{hyperref}       % hyperlinks
\usepackage{url}            % simple URL typesetting
\usepackage{booktabs}       % professional-quality tables
\usepackage{amsfonts}       % blackboard math symbols
\usepackage{nicefrac}       % compact symbols for 1/2, etc.
\usepackage{microtype}      % microtypography
\usepackage{xcolor}     
\usepackage{algorithm}
\usepackage{algpseudocode}
\usepackage{amsmath}
\usepackage{graphicx}
\usepackage{float} % optional, for [H]

%------------------------------------------------------------------------------------------%
\usepackage{xcolor}
\definecolor{dracula}{RGB}{40, 42, 54}
\definecolor{textcolor}{RGB}{235, 235, 235}
\pagecolor{dracula}
\color{textcolor}
%------------------------------------------------------------------------------------------


% Add your own as neccessary

\title{When History Matters: Benchmarking Recurrent RL Agents on Memory Tasks}
% TODO: add your real name and real link
\author{%
  Anh Khoa Pham \\\href{mailto:khoa.pham@uni-hannover.de}{khoa.pham@stud.uni-hannover.de} %\url{https://github.com/khoaphamanh/rl_project}
  \And
  Quang Ky Anh Nguyen  \\\href{mailto:ky.nguyen@stud.uni-hannover.de}{ky.nguyen@stud.uni-hannover.de} %\url{https://github.com/author2}
 % \And
  %Author Three \\ \url{https://github.com/author3}
}

\begin{document}

\maketitle
\section{Motivation}

Standard PPO assumes the Markov property, but many environments are partially observable: the current observation alone is insufficient to act optimally, leading to state aliasing. In \texttt{MiniGrid-Memory}, the agent must recall an object seen earlier to select the correct one at a later fork, a memoryless agent can only guess.

We address this by equipping PPO with a recurrent encoder that maintains a history of past observations. We further ask whether a Transformer encoder handles longer histories better than LSTM, testing this via a reduced observation box ablation.

\textbf{Research Question:} Does equipping PPO with recurrent memory recover performance in partially observable environments, and does replacing LSTM with a Transformer encoder provide an advantage when longer observation histories are required?

\section{Related Topics}

PPO and DQN, introduced in the RL course, operate under the MDP assumption. LSTM and GRU layers can maintain temporal context across time steps, making them natural candidates for POMDPs. Hausknecht et al.\ pioneered this combination with DRQN \cite{drqn}, and Pleines et al.\ (2022) extended it to PPO, introducing RPPO with LSTM/GRU backbones for partially observable environments \cite{rppo}.

\section{Experiments}
\subsection{Enviroment}
We adopt the Memory MiniGrid environment, where the agent begins in a left room with a 7×7 observation field. The task is to identify an object (green key) in the initial room, then proceed to a fork on the right to select the matching object. This instantiates a POMDP: task completion depends on the agent remembering the object's features across time steps. Reward is sparse, provided only upon task success. As shown in Fig. \ref{fig:memory_grid}, we start with \texttt{MiniGrid-MemoryS7-v0} and plan to scale to \texttt{MiniGrid-MemoryS9-v0} or \texttt{MiniGrid-MemoryS11-v0} if preliminary results are favorable.

\begin{figure}[H] % use [htbp] if you don't want to force placement
    \centering
    \includegraphics[width=0.4\linewidth]{memory_grid.png}
    \caption{An example of Memory Minigrid enviroment}
    \label{fig:memory_grid}
\end{figure}

\subsection{Baselines}

We include MLP PPO as a memoryless baseline. Since it conditions only on the current observation without any recurrent state, it cannot resolve state aliasing and can only guess when the task requires recalling past observations.

\subsection{RPPO}

Following Pleines et al., we employ a CNN1D encoder combined with an LSTM layer as the backbone, shared between the actor and critic, with two separate output heads for policy and value. However, we still need to determine through experiments whether sharing the encoder is beneficial or whether splitting it into two independent models yields better performance.

\textbf{Sequence Padding: }
Since episodes have variable lengths until termination, we pad all sequences to a uniform length to satisfy LSTM's input requirements. Embeddings corresponding to padded positions are masked and excluded from loss computation across all three PPO losses (policy loss, value loss, and entropy bonus) \cite{rppo}.

\textbf{Masking Example:}
Given batch size 2:
\begin{align}
\text{Episode 1:} \quad &[t_1, t_2, t_3, \text{PAD}, \text{PAD}] \\
\text{Episode 2:} \quad &[t_1, t_2, t_3, t_4, t_5]
\end{align}

After LSTM encoding and masking:
\begin{align}
\text{Episode 1:} \quad &[0.8, 0.6, 0.4, 0.0, 0.0] \\
\text{Episode 2:} \quad &[0.5, 0.3, 0.7, 0.4, 0.6]
\end{align}

The effective batch contains 8 steps (3 from Episode 1, 5 from Episode 2). The remaining PPO loss terms remain unchanged.

\subsection{Ablation Study}

\subsubsection{Transformer PPO}

We replace the LSTM layer with a fully connected layer followed by a Transformer encoder, keeping all other architectural choices identical to RPPO. Unlike LSTM, the Transformer attends directly over the full observation sequence without compressing history into a fixed-size hidden state.

\subsubsection{Reduced Observation Box (Optional)}

LSTM must compress the entire history into a fixed-size hidden state, which becomes a bottleneck for long sequences. Reducing the observation box from $7\times7$ to $5\times5$ or $3\times3$ narrows the agent's field of view, forcing more steps to solve the environment and thus producing longer sequences. This setting is designed to expose the LSTM bottleneck and test whether the Transformer handles longer dependencies more effectively.


\subsection{Benchmarking}

We benchmark all agents on \texttt{MiniGrid-MemoryS7-v0} with the default $7\times7$ observation box. Table~\ref{tab:bench_main} summarizes our expected outcomes. MLP PPO, being memoryless, cannot resolve state aliasing and is expected to achieve low accumulated return $G_t$. Both LSTM PPO and Transformer PPO maintain a history of past observations and are expected to solve the task, yielding high $G_t$.

\begin{table}[H]
    \centering
    \caption{Expected accumulated return $G_t$ on the default $7\times7$ observation setting.}
    \label{tab:bench_main}
    \begin{tabular}{lc}
        \toprule
        \textbf{Algorithm} & \textbf{Expected $G_t$} \\
        \midrule
        MLP PPO       & Low  \\
        LSTM PPO      & High \\
        Transformer PPO & High \\
        \bottomrule
    \end{tabular}
\end{table}

Optionally, if time permits, we evaluate on a reduced $5\times5$ observation box to stress-test the memory mechanisms under longer sequences. We expect LSTM PPO to struggle due to its fixed-size hidden state bottleneck, while Transformer PPO should maintain high $G_t$ thanks to its attention over the full sequence. Table~\ref{tab:bench_ablation} summarizes the expected outcomes.

\begin{table}[H]
    \centering
    \caption{(Optional) Expected accumulated return $G_t$ on the reduced $5\times5$ observation setting.}
    \label{tab:bench_ablation}
    \begin{tabular}{lc}
        \toprule
        \textbf{Algorithm} & \textbf{Expected $G_t$} \\
        \midrule
        LSTM PPO $5\times5$        & Low  \\
        Transformer PPO $5\times5$ & High \\
        \bottomrule
    \end{tabular}
\end{table}

\section{Timeline}

Our project timeline is as follows:
\begin{itemize}
    \item Literature review: 3 days
    \item Implementation: 4 days
    \item Runing code: 1 week
    \item Reporting, visualization, and poster design: 2 days 
\end{itemize}

\bibliographystyle{IEEEtran}
\bibliography{references}

\end{document}
